% Pursuing Stacks (cotes 134-2 to 134-8): archive cote 134-5, site page 132 (PDF page 133), typescript page 319.
% typescript. Draft compared: ps5.tex:1327-1377.
% AI-transcribed by Claude Opus 5.5 (claude-opus-5-5, adaptive thinking, effort high), Message Batches API, batch ps3 (msgbatch_01XdFmuVVBAcgc7EFxdtCqsj), 2026-11-12; reviewed against both supplied scan halves; see CHECK_DECISIONS.tsv.
% The hemispherical sign has a double oval outline; bigcirc is the supplied glyph approximation.
\scanpage{134-5/132}
\mspage{319}

too that I never got around introducing formall\over{y}{y} the hemispherical test
category, which presumably will be \over[typed]{b}{v}\unsure{a|e}ry useful when it comes to study-%
ing stacks - this too could be done in part IV, as well as proving
of course that \hand{$\bigcirc$} is a strict test category indeed, or better still,
a contractor. It may be fun too constructing test functors from any
one of the three basic test categories in the trinity, to the category
of complexes defined by the two others - six cases altogether to con-%
sider! But as I am not in the process of writing the ``Elements d'Algè-%
bre Homotopique'' , maybe I will skip this !

In the same section 43, I raised the question as to whether the
ordered set of all non-empty finite subsets of a given infinite set,
viewed as \over[typed]{g}{a} category in the usual way \over{f}{f}or an ordered set, was a weak
test category (on page 78 it was seen not to be a test category, and
it is immediate then that it is not totally aspheric either). O\(\underset{\text{n}}{\text{n}}\)e in-%
teresting application, as noticed there, would be to the effect that
(Ord), the subcategory of (Cat) defined by ordered sets, is a modeli\deleted{z\unsure{a}}%
zer (for the induced notion of weak equivalences). \xout{Here}\inserted[typed]{Now} the question
arises moreover whether this modelizing structure comes from an aspheri-%
city structure, or even from a contractibility structure - and the
same questio\over[typed]{j}{n} arises in the more general situation described in the
proposition of page 74.

A last question \over[typed]{I}{a}long these lines I would like to clear up,
is the relation of total \inserted[typed]{$\underline{W}$-}asphericity for an asphericity structure,
for variable $\underline{W}$ \(\underset{\hand{\smile}}{\text{,}}\) when $\underline{W} \hand{\subset} \underline{W}'$ . Assuming the localizers satisfy
the condition L\(\underset{\hand{\illegible{1}}}{\text{o}}\)c 4 ), is it true that \inserted[typed]{total} $\underline{W}$-asphericity is equivalent
to total $\underline{W}'$-asphericity - \xout{and} or equivalently, is it equivalent
to total 0-connectedness ?

\margin{132}
